% Source: https://artofproblemsolving.com/community/c3h1927430_problem_need_help
% Stable identity: aops-topic-1927430-q01
Rewrite the equation as $\dfrac{1}{x}+\dfrac{2}{y}=1+\dfrac{3}{z}$, so $\dfrac{1}{x}+\dfrac{2}{y}>1$.

\textbf{Case $x=1$.} The equation becomes $\dfrac{2}{y}=\dfrac{3}{z}$, i.e.\ $3y=2z$. So $y$ is even, $y=2k$, and then $z=3k$. This gives $(1,2k,3k)$ for every positive integer $k$.

\textbf{Case $x\ge2$.} Then $\dfrac{2}{y}>1-\dfrac{1}{x}\ge\dfrac{1}{2}$, so $y<4$.
\begin{itemize}
\item $y=1$: $\dfrac{1}{x}+2-\dfrac{3}{z}=1$ gives $\dfrac{3}{z}=\dfrac{x+1}{x}$, so $z=\dfrac{3x}{x+1}$. Since $\gcd(x,x+1)=1$, we need $x+1\mid3$, so $x=2$ and $z=2$: the triple $(2,1,2)$.
\item $y=2$: $\dfrac{1}{x}+1-\dfrac{3}{z}=1$ gives $z=3x$: the triples $(k,2,3k)$.
\item $y=3$: $\dfrac{1}{x}+\dfrac{2}{3}-\dfrac{3}{z}=1$ gives $\dfrac{3}{z}=\dfrac{1}{x}-\dfrac{1}{3}$, which must be positive, so $x<3$, i.e.\ $x=2$. Then $\dfrac{3}{z}=\dfrac{1}{6}$ and $z=18$: the triple $(2,3,18)$.
\end{itemize}

So the solutions are
\[
\boxed{(1,2k,3k)\ \text{and}\ (k,2,3k)\ \text{for positive integers }k,\quad\text{together with }(2,1,2)\text{ and }(2,3,18).}
\]
Check: $\dfrac{1}{2}+2-\dfrac{3}{2}=1$ and $\dfrac{1}{2}+\dfrac{2}{3}-\dfrac{1}{6}=1$; the two families satisfy the equation identically.
